% Transcription — fonds Alexandre Grothendieck, shelfmark 7, batch 3 (pages 41-60).
% Established from the facsimile archives/batches/7/batch-03.pdf, cut from a
% copy of 7.pdf fetched through the project's own relay
% (https://grothendieck-relay.onrender.com/source/7.pdf) rather than from
% grothendieck.umontpellier.fr directly; per the skill transcribe-grothendieck.
% The batch file carries no cover sheet: PDF page k is the archivists' page
% 40+k, and their pencil numbers bottom-left agree throughout. The folder is
% seventy-one pages; this is the third of four batches (the others are
% transcribed in parallel by other passes).
%
% Pass: Opus 5.5 (claude-opus-5-5), 2026-09-26 — first pass, unchecked against the pages by a human.
%
% The folder sits in the inventory group "4-9" « Cristaux »; it has its own
% date in the catalogue, which \dating{} carries verbatim, with no group
% sentence.
%
% Pages skipped: none. Pages summarised: none. All twenty leaves are his own
% manuscript in blue ink and are transcribed in full, as far as the hand
% allows. The bottom halves of pp. 49 and 55 are blank.
%
% His pagination and dates. No page number of his on any leaf. Dates:
%   p. 56 carries at the top right « 6.12.1966 » (the third digit poorly
%   formed), under the heading « Laïus »; it agrees with batch 2's
%   « ?5.11.66 » on p. 35. \dating{} stays the catalogue's.
% Runs, by his own headings:
%   41-43  continuation of batch 2's « L'idée de W.M. » (p. 41 opens by
%          finishing the sentence « (les idéaux sont-ils automatiquement
%          fermés ?) » broken off at the foot of p. 40): the W-M lifting
%          theorem a)-d) for dagger algebras, its corollary, W-M's local
%          and global invariants; « 3) Critique de W.M. » (p. 43-46).
%   46-49  « §3 [struck 4)] Relations heuristiques de W.M. avec une idée de
%          Gauss-Manin », points 1-4; p. 49 lists 2-4 as headings only.
%   50-52  point 3.5 (Manin, the Cartan homotopy formula, the Gauss-Manin
%          connection on R^i f_*(Omega_{X/S})), remarks, and the summing-up
%          « puzzle » a)-c) of p. 52.
%   53-55  « §4 Topos stratifiant et topos cristallin », 4.1-4.3: the
%          stratification as a system of extensions to infinitesimal
%          neighbourhoods, the crystalline and stratifying sites, the theorem
%          H_strat = H_DR in char. 0 and the conjecture H_strat = H(X^an, C).
%   56-58  « Laïus », 6.12.1966: stratifying cohomology as Zariski
%          hypercohomology of the cosimplicial C^*(F); the pro-Module Q(F)
%          functorial in differential operators; P^i, Q^i[M].
%   59-60  a fresh plan §1-§8 of the whole series of talks (§5 « Site
%          cristallin et site stratifiant », §6 « Calculs Čechistes », §7
%          comparison De Rham / crystalline, §8 back to char. p), then 6.1
%          « Coh. stratifiante »: the Leray spectral sequence of the covering
%          object X~, its terms computed by Mittag-Leffler, and the boxed
%          H*(X_strat, F) = H*(nu -> F(X^(nu+1)/X)) for X affine. P. 60
%          breaks off mid-sentence (« on a une ») into batch 4.
%
% What the batch is about. The end of the Washnitzer-Monsky discussion
% (overconvergent W^dagger-algebras, unique liftings up to non-unique
% isomorphism, invariance of H(Omega(A/W)) mod torsion; the site of W-M
% liftings of special affine opens) and his critique of it (W-coefficients
% and a p-adic torsion, but no Leray spectral sequences, no finiteness, no
% duality or Künneth); the heuristic link with Gauss-Manin (an invariance
% under change of lifting read as a stratification; Manin's connection,
% true in char. 0, false in char. p for elliptic families); then the first
% statement of the crystalline and stratifying sites and topoi, the
% comparison theorem with De Rham in char. 0 and the conjecture for
% singular X over C, and the Čech computation of stratifying cohomology.
%
% Notation fixed once for the batch (continuing batch 2): W^{\dagger} for
% the overconvergent ring; \mathfrak{X} for his formal scheme; \mathbb{H}
% for hypercohomology; \Omega^{\cdot} for the dotted Omega; \mathcal{K},
% \mathcal{H}^i for his script K and H (p. 47-48); \underline{E} for his
% underlined E; \mathcal{P}^i, \mathcal{Q}^i for his round P and Q (p.
% 58); X_strat, X_zar for his subscripts; \Delta^{(\nu)}_X(i) for the i-th
% infinitesimal neighbourhood of the diagonal.
%
% Hardest pages and readings to check first: the interlinear insertions of
% p. 43 and 44; the struck lines of p. 48 and 54 (the cancelled R f_*
% formulas, boxed and hatched); the vertical margin note of p. 58 (only
% « Anti-Riesz » reads); the dense bottom of p. 57; the superscript sign
% « (Sigma) » on p. 47; the glosses under the formulas of p. 59.
% Apparatus: 135 \ill{} and 71 \uncertain{}.

\documentclass[11pt,a4paper]{article}
% Le préambule commun à toutes les transcriptions du fonds Grothendieck.
%
% Il tient en une idée : **l'incertitude se compose, elle ne se résout pas.**
% Une transcription qui lisse un mot douteux a détruit la seule chose pour
% laquelle on la faisait — un lecteur doit pouvoir savoir, à chaque endroit,
% ce qui était sur le feuillet et ce qui a été ajouté. Les six macros qui
% suivent sont donc l'appareil critique, et non de la décoration.
%
% Les mêmes commandes sont reconnues par `scripts/render.mjs`, qui produit la
% vue de lecture du site. Une macro ajoutée ici doit l'être aussi là-bas,
% faute de quoi le rendu échouera bruyamment — ce qui est voulu : un rendu
% silencieusement faux est pire qu'un rendu absent.

\ProvidesPackage{grothendieck}[2026/08/08 Transcriptions du fonds Grothendieck]

\usepackage{amsmath,amssymb,amsthm}
\usepackage{mathrsfs}
\usepackage{stmaryrd}
% Les diagrammes commutatifs sont partout dans ce fonds : c'est la langue
% même de la géométrie algébrique de Grothendieck, et une transcription qui
% les rendrait en prose ne transcrirait rien.
\usepackage{tikz-cd}
% Français par défaut — c'est la langue des feuillets — mais l'anglais est
% chargé aussi : la traduction et le résumé sont des documents anglais, et la
% typographie française (espaces avant les deux-points) y serait une faute.
% Ces documents basculent par \selectlanguage{english} en tête de corps.
\usepackage[english,french]{babel}
\usepackage{fontspec}
\usepackage{geometry}
\geometry{a4paper,margin=25mm,marginparwidth=28mm}
\usepackage{marginnote}
\usepackage{xcolor}
% ulem, not soul. `soul`'s \st and \ul are built on a character-by-character
% scan that XeLaTeX's font machinery breaks: the document fails with an
% undefined control sequence rather than degrading. ulem does the same two jobs
% and is Unicode-engine safe. `normalem` stops it hijacking \emph, which the
% transcriptions use for Grothendieck's own emphasis.
\usepackage[normalem]{ulem}
\usepackage{hyperref}
\hypersetup{colorlinks=true,linkcolor=black,urlcolor=black,pdfborder={0 0 0}}

% --- Métadonnées du lot -----------------------------------------------------
% Reprises telles quelles par le script de rendu, qui les affiche en tête de
% la vue de lecture. Elles fixent ce que le fichier prétend couvrir : sans
% elles, une transcription flotte sans point d'attache dans un fonds de seize
% mille feuillets.

\newcommand{\folder}[1]{\gdef\grothFolder{#1}}
\newcommand{\batch}[1]{\gdef\grothBatch{#1}}
\newcommand{\leaves}[2]{\gdef\grothFirst{#1}\gdef\grothLast{#2}}
\newcommand{\foldertitle}[1]{\gdef\grothFoldertitle{#1}}
\gdef\grothFolder{?}\gdef\grothBatch{?}\gdef\grothFirst{?}\gdef\grothLast{?}\gdef\grothFoldertitle{}

% --- Le résumé --------------------------------------------------------------
%
% Il ouvre la lecture modernisée, sous le titre « Résumé », et il est composé
% en petit. Ce n'est pas un ornement : c'est le seul passage du document qui
% ne soit pas des mathématiques, et un lecteur doit pouvoir le reconnaître
% avant de l'avoir lu — puis le sauter s'il connaît déjà le sujet, ou s'y
% arrêter s'il ne le connaît pas. Le filet qui le suit marque la reprise du
% corps.

\newenvironment{resume}
  {\small\setlength{\parskip}{0.45em}}
  {\par\medskip\hrule\medskip}

% --- Mots-clés ---------------------------------------------------------------
%
% En anglais, en clôture du résumé de la lecture modernisée : le vocabulaire
% moderne sous lequel chercher ce que les feuillets construisent. Le manifeste
% du site (`npm run manifest`) les extrait pour étiqueter la cote entière —
% cette ligne est la source unique des tags, il n'y a pas de fichier à côté.
\newcommand{\keywords}[1]{\par\smallskip\noindent
  {\footnotesize\textsc{Keywords} --- \textit{#1}}\par}

% --- L'appareil critique ----------------------------------------------------

% Le numéro de feuillet, en marge. C'est l'ancre qui permet de régler un
% désaccord sur une formule en regardant *un* feuillet plutôt qu'une chemise
% de six cents. Le site s'en sert aussi pour tourner les pages du fac-similé
% quand on fait défiler la transcription.
\newcommand{\leaf}[1]{%
  \marginnote{\footnotesize\color{gray!70}\textsf{#1}}%
  \label{leaf:#1}\ignorespaces}

% Illisible : on ne devine pas. Un mot inventé qui se lit comme les autres est
% le pire résultat possible.
\newcommand{\ill}{\textcolor{red!60!black}{[\,\ldots\,]}}

% Lecture douteuse : le mot est donné, et signalé comme incertain.
\newcommand{\uncertain}[1]{\uline{#1}}

% Ajout de l'éditeur — une lettre manquante, un mot sous-entendu.
\newcommand{\add}[1]{\textcolor{blue!50!black}{[#1]}}

% Biffé par Grothendieck lui-même. Ce qu'il a rayé fait partie du document :
% une rature dit souvent où la pensée a changé de direction.
\newcommand{\struck}[1]{\sout{#1}}

% Note du transcripteur, sur l'état matériel du feuillet.
\newcommand{\note}[1]{\par\smallskip{\small\color{gray!60!black}\textit{[#1]}}\par\smallskip}

% Note marginale de Grothendieck, distincte du corps du texte.
\newcommand{\marginal}[1]{\par\smallskip\noindent\hspace{1em}%
  {\small\color{gray!60!black}#1}\par\smallskip}

% --- Titre ------------------------------------------------------------------

\AtBeginDocument{%
  \begin{center}
    {\large\bfseries Fonds Alexandre Grothendieck --- cote n\textsuperscript{o}~\grothFolder}\\[2pt]
    {\itshape\grothFoldertitle}\\[4pt]
    {\small feuillets \grothFirst--\grothLast\ (lot \grothBatch)}\\[2pt]
    {\footnotesize\color{gray!60!black}Transcription établie d'après le fac-similé numérisé
     par l'Université de Montpellier.\\
     Les lectures incertaines sont soulignées, les illisibles notées [\,\ldots\,],
     les ajouts entre crochets.}
  \end{center}
  \vspace{1em}}

\endinput


\folder{7}
\batch{3}
\pages{41}{60}
\dating{1966-1967}
\watermark{Édition de démonstration}
\foldertitle{Notes laïus novembre-décembre 1967 (Cristaux) : notes manuscrites (1966-1967)}

\begin{document}
\page{41}

\uncertain{automatiquement} fermés\,?).\note{suite de la phrase
interrompue en bas de la p.\ 40 (batch 2) : « (les idéaux sont-ils
automatiquement fermés\,?) »} \struck{W.M. prouvent}

Il semble que le théorème suivant soit vrai (dont W.M. prouvent une
variante plus faible) :

a) Si $A_0$ de t.f.\ sur $k$ et lisse sur $k$,
$\exists$ $A$ \add{W-M}\note{ajouté au-dessus de la ligne} algèbre sur $W$,
\emph{plats} sur $W$, et un $k$-isomorphisme
\[
  \varphi : A \otimes_W k \simeq A_0 .
\]

b) De plus, un tel couple $(A, \varphi)$ est unique (à isom.\ (non
unique) près), i.e.\ si on a deux $(A', \varphi')$, il existe un
$W$-isom.\ $u : A \to A'$ rendant commutatif le diagramme
\uncertain{évident}.

c) Définissons, pour une algèbre de W.M.\ $A$ sur $W$,
\add{($\Omega^{1}(A/W)$ par la propriété} universelle pour les
dérivations de $A/W$ dans les $A$-modules de type fini, d'où
\[
  \Omega^{*}(A/W) = \textstyle\bigwedge^{*}(A/W) ,
\]
\note{ainsi sur la page : le $\bigwedge$ porte l'étoile en exposant, et
l'on attendrait $\bigwedge^{*}\Omega^{1}(A/W)$} qui est un \emph{complexe}
de $W$-algèbres.

\struck{b) \ill{} \ill{} divisions \ill{}}\ \struck{Soient $(A, \varphi)$
$(A', \varphi')$ \ill{} \ill{} \ill{} des}\note{deux lignes barrées ;
au-dessus de la seconde, « Soient » et « $A$. $A'$ », eux aussi barrés}
\struck{\ill{} : $A \to A'$ \ill{} \ill{}}

\struck{\ill{}} $u : A \to A'$ un homom.\ de W.M.-\uncertain{d)}\note{la
lettre, dans la marge gauche, se lit « b) » ou « d) »} algèbres lisses sur
$W$, d'où
\[
  u^{*} : H^{\cdot}\bigl(\Omega^{\cdot}(A'/W)\bigr) \longrightarrow
  H^{\cdot}\bigl(\Omega^{\cdot}(A/W)\bigr)
\]
\note{la flèche est un simple trait, sans pointe ; l'ordre $A'$, puis $A$,
est celui de la page}

\page{42}

Cet homomorphisme ne dépend que de l'homomorphisme
$u_0 : A_0 \to A'_0$, induit sur les \uncertain{algèbres réduites} mod $p$.

\emph{Cor.} Si $A$ est un relèvement de W.M.\ de
\add{\uncertain{lisse de}}\note{dans la marge gauche, deux mots reliés par
un trait} $A_0$ lisse sur $k$, \uncertain{alors}
$H^{\cdot}\bigl(\Omega^{\cdot}(A/W)\bigr)$ ne dépend, à isom.\ canonique
près, que de $A_0$.

En fait, W-M.\ ne prouvent ces résultats que pour des $A_0$
« spéciales », --- mais il y a \uncertain{assez} de celles-ci. De plus,
dans leurs \uncertain{notes}, ils ne prouvent la propriété d'invariance
\ill{} que mod torsion, i.e.\ en tensorisant par $K$.

\uncertain{(4)}\note{signe peu lisible en tête de l'alinéa, dans la marge}
Lorsque $X_0$ est un schéma lisse quelconque sur $k$, en utilisant ces
résultats, W.M.\ définissent \add{dans [~]}\note{au-dessus de la ligne,
crochets laissés vides pour une référence, et relié par un trait} des
« faisceaux de cohom. » \struck{$H^{*}(X$} $\mathcal{H}^{*}_{X_0}$ sur
$X_0$, qui sont des faisceaux de $W$-modules.

Cela ne donne pas encore d'invariants

\page{43}

globaux, comme on va \uncertain{le} voir en avançant.
\struck{Dernièrement} \add{Récemment} W.M.\ ont utilisé la notion de
\struck{topos} site pour définir des invariants globaux : \emph{le site
des relèvements de W.M.\ des ouverts affines spéciaux}, pour la
« topologie de Zariski ». Sur ce site, les
$\Omega^{\cdot}_{A/W}$\note{l'indice se lit « $A/k$ » ou « $A/W$ »}
forment un complexe de faisceaux, et ils \uncertain{prennent}
l'hypercohomologie. Ils trouvent des invariants globaux
\ill{}, qui sont des modules sur $W$. Ils \add{\uncertain{vérifient} (?)}
que dans le cas envisagé au début, \add{\uncertain{ils retrouvent bien}}
$\mathbb{H}^{*}(\mathfrak{X}, \Omega^{\cdot}_{\mathfrak{X}/S})$, d'où la
\add{propriété d'invariance} W.M.\ que j'y avais énoncée.\note{les deux
ajouts sont en interligne, serrés ; l'ordre des mots est reconstitué}

3) \emph{Critique de W.M.}

C'est une idée nouvelle importante, en ce qu'elle met en relation la
coh.\ de D.R.\ en car.\ $0$, avec la coh.\ \uncertain{algébrique}
$p$-adique hypothétique. Elle fournit une théorie cohomologique
« $p$-adique en car.\ $p$ », qui évite une des critiques principales
\uncertain{qu'on avait faites} à D.R.\ naïve : elle

\page{44}

a comme anneau de coeff.\ $W(k)$, dont le corps des fractions $K$ est de
car.\ $0$, et non $k$. De plus, \uncertain{ces} $W$-modules
\struck{ne} \add{(s'il est \uncertain{vrai qu'elle vit} sur $W(k)$, pas
seulement $K$)}\note{en interligne} \struck{sont} fournit
\struck{des relations satisfaisantes} une notion de torsion $p$-adique,
qui est peut-être \struck{\ill{} la \ill{}} la bonne : est-ce la torsion
$p$-adique de la fibre générique\,? géométrique\,? \ill{}
\ill{}\note{deux mots, en bout de ligne et au début de la suivante}
question, extrêmement importante et profonde, des relations entre la
cohomologie de D.R.\ et la cohomologie $p$-adique, \ill{} \ill{}
spécialisation de la car.\ $0$ à la car.\ $p > 0$. Cette question déjà
\add{dans un \uncertain{langage} différent,}\note{en interligne, au-dessus
de « l'an\ldots »} envisagée par Tate, (dans \ill{} \ill{} l'an\add{née}
dernière, dans la \ill{} tout au moins).\note{le mot après « dans la » est
entouré d'une courbe ; on y lit deux capitales, « HH » ou « IH »}

Sur l'importance de cette question \struck{\ill{}} pour \ill{} de
certains pbs liés aux \add{\emph{familles de}}\note{au-dessus de la ligne,
relié par un trait} cycles algébriques, cf.\ \ill{} \ill{} aux Invariants
\add{[~]}\note{crochets laissés vides, au-dessus de la ligne, reliés par un
trait oblique} \ill{} parfaits, prochainement.

\page{45}

Il reste cependant des \struck{graves} critiques : formelles, la plus
importante étant \ill{} \ill{} \ill{}\note{fin de la ligne, sans doute
« de voir »} fournie par le différentielle, ce qui pratiquement limite la
théorie aux $X_0$ lisses. Ne \struck{\ill{}} travaillant qu'avec les seuls
complexes $\Omega^{\cdot}_{X_k}$ de \struck{différentielles} De Rham,
\add{il est sans espoir que}\note{en interligne, au-dessus de la ligne
suivante, et relié à celle-ci} \struck{cette} « W-\ill{} »
\struck{permet pas d'écrire} \add{fournisse} des suites spectrales
\struck{De Rham} type Leray pour un morphisme $f_0 : X_0 \to Y_0$, --- si
précieuses pour faire des \uncertain{dévissages} dans les démonstrations.

Comme conséquences, de \struck{cela} ceci, sans doute, W.M.\ n'ont pas été
capables de prouver pour leur cohomologie les \struck{\ill{}}
propriétés formelles les plus \uncertain{élémentaires}, \ill{} \ill{}
\ill{} d'entre elles est certainement le fait que leurs
$H^{i}(X)_{\mathrm{WM}}$ sont des modules de t.f.\ sur $W$ --- donc sans
doute

\page{46}

non plus le théorème de dualité, ni de Künneth. \struck{Je \ill{}
d'\ill{} (}

(Il semble cependant \ill{} par un \ill{}, ils envisagent d'utiliser leur
cohomologie pour donner une expression de la \uncertain{formule}
Lefschetz-Weil de la fonction $\zeta$ d'un \uncertain{schéma} projectif
et lisse sur un corps fini, en \ill{} que les Frobenius opèrent sur les
$H^{i}(X)$ \ill{} des opérateurs : \ill{} \ill{}\ldots\,! Ce formalisme
semble \ill{} \ill{} de celui de Dwork, que je ne connais pas\ldots).

\section*{§\,3 \struck{4)} Relations \add{heuristiques de W.M.} avec une
idée de Gauss-Manin.}
\note{« heuristiques de W.M. » est écrit au-dessus et inséré par un
trait entre « Relations » et « avec » ; le titre est souligné. Le « 4) »
biffé est récrit en « §\,3 »}

1. Caractère assez surprenant à première vue de l'assertion d'invariance
des W.M. \uncertain{Supposons} p.\,ex.\ que l'on ait

\page{47}

une grosse variation \ill{} de $X_0$, s'exprimant \struck{disant} par le
fait que $X_0$ est la fibre, en un point $t_0$ au dessus du pt fermé $s_0$
de $S$, d'un schéma propre et lisse $Y$ \add{/}\note{un trait oblique
relie « $Y$ », récrit au-dessus, à « sur »} sur un $M$ de type fini sur
$S$. Les différentes sections de $M$ sur $S$ \uncertain{donnent}
différentes façons, de relever $X_0$ en un $X$ propre et lisse sur $S$.
Prenant
\[
  (*) \qquad \mathcal{K}^{\cdot} = Rf_{*}(\Omega^{\cdot}_{Y/M})
\]
on voit que pour toute telle section $g$ de $M$ sur $S$, définissant
\struck{un} $X_g^{(\Sigma)}$, et
\[
  f_g^{(\Sigma)} : X_g^{(\Sigma)} \longrightarrow S
\]
\note{l'exposant de $X_g$ et de $f_g$ est un petit signe entre
parenthèses, lu $\Sigma$ sous réserve ; sur $f_g$ il est d'abord écrit,
biffé, puis récrit}
on a
\[
  Rf_{g*}(\Omega^{\cdot}_{X_g/S}) \xleftarrow{\ \simeq\ } g^{*}(\mathcal{K}^{\cdot})
\]
\note{sous la flèche : « isom.\ can. »}
($g^{*}$ au sens des catégories dérivées\ldots). Il s'ensuit, du moins
modulo torsion, que
\[
  R^{i}f_{g*}(\Omega^{\cdot}_{X_g/S}) \xleftarrow{\ \simeq\ } g^{*}(\mathcal{H}^{i})
\]

\begin{tikzcd}[column sep=large]
Y \arrow[r, "f"] & M \arrow[r, "h"] & S = \operatorname{Spec} W
\end{tikzcd}

\note{croquis dans la moitié gauche de la page : un cube marqué $Y$, une
flèche $f$ vers un carré $M$ traversé d'un trait vertical et d'un faisceau
de droites passant par un point, puis une flèche verticale $h$ vers un
segment pointé, marqué « Spec $W = S$ » ; redessiné ici en ligne. En haut
à droite, au crayon très pâle, un petit calcul sans rapport apparent
(« $(k+1)$ \ldots{} $\frac{k}{2}$ »), illisible}

\page{48}

\emph{Ainsi, on trouve que}, \add{(\emph{posant} $(*)$),} \emph{les
complexes $g^{*}(\mathcal{K}^{\cdot})$, pour les différentes sections $g$
de $M$ sur $S$ passant par $t_0$, sont canon.\ isomorphes entre eux (dans
une des cat.\ dérivées), donc les $g^{*}(\mathcal{H}^{i})$}
(où $\mathcal{H}^{i} = \mathcal{H}^{i}(\mathcal{K}^{\cdot})$
\add{$= R^{i}f_{*}(\Omega^{\cdot}_{Y/M})$})\note{l'égalité
$= R^{i}f_{*}(\Omega^{\cdot}_{Y/M})$ est ajoutée au-dessus, en bout de
ligne, et reliée par un trait} \emph{sont canoniquement isomorphes entre
eux}, \struck{pour} \emph{mod torsion}.\note{toute la phrase est
soulignée, ligne à ligne} La donnée d'un tel système transitif
d'isomorphismes \struck{\ill{}} peut être considérée comme une
\emph{structure} \ill{} sur
\[
  \mathcal{K}^{\cdot} = Rf_{*}(\Omega^{\cdot}_{Y/M})
\]
[en \ill{} \uncertain{relativement}, à la base $S$], resp.\ sur les
$\mathcal{H}^{i} = R^{i}f_{*}(\Omega^{\cdot}_{Y/M})$,
\struck{Une telle} \struck{\ill{} \ill{} \ill{} p.\,ex.\ \ill{}}
\add{\ill{} \ill{}}\note{ligne biffée, avec un ajout au-dessus en partie
illisible} à une \emph{structure de descente}. Une structure de
\struck{\ill{}} \ill{} \ill{} \add{sur un $M$-module $\underline{E}$
disons,}\note{en interligne} type, \ill{} \ill{} \ill{} qui \ill{} ses
fibres \ill{}, i.e.\ définissons un isom.
\[
  \underline{E} \simeq h^{*}(\underline{E}_0)
\]
\note{après $\underline{E}$, un $\mathcal{H}$ biffé}
($\underline{E}_0$ un $M_0$-module \ill{} sur $S$], exprime un

\page{49}

tel système transitif d'isomorphismes entre les $g^{*}(\underline{E})$.
Nous nous \ill{} décrire des structures, \add{\uncertain{plus}} fortes que
celles ci-dessus, qui définissent encore un tel syst.\ transitif,
disons : les \emph{structures de stratification relativement à} $S$.

2. Définition générale d'une stratification. \add{Relations avec
connexions.}\note{en interligne, d'une encre plus pâle}

3. Système transitif d'isom.\ qu'elle fournit, si \ill{} est un \ill{}
\ill{} local complet.

4. La \struck{idée de} connexion de Gauss-Manin.\note{le « 4 » est récrit
sur un autre chiffre ; le bas de la page est blanc}

\page{50}\note{la page commence au milieu d'un développement, sans titre ;
le « 3.5 » qu'elle porte plus bas reprend la numérotation des points 1--4
annoncés p.\ 49, dont le début (l'exposé de Manin) ne figure pas sur les
feuillets précédents}

Dans le cas où $k$ est de car.\ $0$, cela se donne une connexion
intégrable absolue sur $\mathcal{H}^{*}$ (Manin n'utilise pas ce langage
géométrique). Puis, dans le cas des $H^{1}$, et lorsque \struck{$K$ est}
provient d'une variété \add{propre}\note{au-dessus de la ligne} lisse, Manin
montre de façon assez compliquée que
$\mathbb{H}(X, \Omega^{\cdot}_{X/k})$ (qui est
$\subset H_{\mathrm{DR}}(K/k)$) est stable par les opérations
$\theta_X$.\note{la lettre sous $\Omega$, $h$ ou $k$, n'est pas nette ;
de même le « $K$ » de $H_{\mathrm{DR}}(K/k)$}

3.5. En fait, par des arguments très simples directs, \struck{que}
reposant sur la « formule d'homotopie » de Cartan
\[
  d\, i_X + i_X\, d = \theta_X
\]
on peut montrer, pour un morphisme $f : X \to S$ :

(i) Pour toute dérivation $\xi$ de $\underline{\mathcal{O}}_S$ qui se
remonte, \add{loc.\ au dessus} \struck{\ill{}} de $S$, en une dérivation
$\overline{\xi}$ de $\underline{\mathcal{O}}_X$ [condition toujours
satisfaite si $X$ est \struck{\ill{}} \emph{formellement lisse sur $S$} et
\emph{affine} sur $S$] l'effet de $(\overline{\xi}, \xi)$ opérant sur
$f_{*}(\Omega^{\cdot}_{X/S})$\note{le symbole devant $f_{*}$ est noirci par
une rature}, \struck{est} indépendant \emph{à une homotopie} \struck{\ill{}}
\emph{près} du choix de $\overline{\xi}$, \add{dans différent
(\uncertain{si} $f$ \uncertain{formellement} lisse et affine)}\note{ajout
entouré, en interligne, en partie douteux}, dépendant de $\xi$ ; donc en
particulier, les
\[
  \mathcal{H}^{i}\bigl(f_{*}(\Omega^{\cdot}_{X/S})\bigr) =
  R^{i}f_{*}(\Omega^{\cdot}_{X/S})
\]
sont des Modules sur lesquels les dérivations de $S$ opèrent, --- et on
vérifie la loi de linéarité et de crochet trivialement. Si $S$ lisse sur un
$T$, cela donne pour $f_{*}(\Omega^{\cdot}_{X/S})$ et
\uncertain{via} les $R^{i}f_{*}(\Omega^{\cdot}_{X/S})$ une conn.\
intégrable rel/$T$.\note{les deux dernières lignes sont serrées en bas de
page : « crochet trivialement » et « donne pour » sont écrits l'un sous
l'autre}

\page{51}

(ii) Par une construction reposant sur le \uncertain{même} principe, mais
seulement un peu plus sophistiquée techniquement, (cf.\ appendice), on
trouve que si $f$ est lisse \uncertain{propre} et qu.-sép., les
\[
  Rf_{*}(\Omega^{\cdot}_{X/S})
\]
est muni d'une \emph{connexion} \emph{absolue canonique} [au sens des
catégories dérivées]. Cela ne signifie pas tout à fait que les
$R^{i}f_{*}(\Omega^{\cdot}_{X/S})$ sont munis individuellement de
connexions absolues ; c'est pourtant vrai si $S$ de car.\ $0$, et $f$
propre et lisse (car alors les formations des
$R^{i}f_{*}(\Omega^{\cdot}_{X/S})$ eux-mêmes commute à l'extension de la
base\ldots) --- \ill{} \ill{}, bien sûr, dans le cas lisse affine envisagé
en (i)\ldots{} La vérification (pour $f$ propre, lisse) que cette
connexion est compatible avec la construction transcendante, est en fait
triviale.

\emph{Remarques.} On pourrait espérer que ces connexions canoniques
proviennent en fait de \emph{stratifications} canoniques, de façon que
l'on aurait une

\page{52}

explication « à priori » de la \struck{construction de}
\add{invariance}\note{au-dessus du mot biffé} affirmée par
Washnitzer-Monsky. C'est bien vrai si $S$ est de car.\ zéro, mais
malheureusement c'est faux en car.\ $p > 0$, --- si on considère par ex.\
des projections structurales de \add{familles de}\note{en interligne,
relié par un trait} courbes elliptiques (\uncertain{schémas} elliptiques)\,!
Il s'agit donc pour le moment d'une relation heuristique et
sentimentale, entre deux idées différentes, dont les relations exactes
\struck{restant} \struck{\ill{}} \add{\uncertain{laissent}}\note{au-dessus
de la ligne} encore \add{d'avoir}\note{au-dessus de « été », relié par un
trait} été bien élucidées\ldots\ldots


Pour résumer, on se trouve conduit devant un puzzle, où il s'agit de
reconstituer une théorie d'ensemble \struck{\ill{} \ill{}} \struck{et}
combinant

a) Une théorie cohomologique des complexes d'opérateurs différentiels,
\struck{qui} (\uncertain{notamment} en vue d'\uncertain{extension} au
complexe de D.R.)\note{la parenthèse est en interligne}

b) Une étude systématique des structures stratifiées, notamment
\add{\uncertain{en vue}}\note{au-dessus de « \ldots{}les »}
\struck{\uncertain{souples}} qui \add{(\uncertain{recevront des}
explications conceptuelles \struck{\ill{}})}\note{ajout entouré, en
interligne, relié par un trait} \ill{} \ill{}

c) La théorie de Washnitzer-Monsky.

\page{53}

\section*{§\,4 \emph{Topos stratifiant} et \emph{topos cristallin}.}

4.1. Interprétation d'une stratification relative d'un Module
$\underline{F}$ sur $X/S$. Si \ill{} \ill{} une telle stratification,
alors pour tt $S$-morphisme $f : Y \to X$, et tt voisinage
infinitésimal $Y'$ de $Y$, \add{(sur $S$)}\note{au-dessus de la ligne}
\struck{qui \uncertain{puisse}} tels que $f$ \uncertain{veuille} bien se
prolonger (localement au moins) en $f' : Y' \to X$, un système
\emph{transitif} d'isom.\ entre les $f'^{*}(\underline{F})$. On peut donc
dire que la \struck{\ill{}} \add{\uncertain{donnée} d'une}\note{en
interligne, au-dessus du mot biffé} stratification \add{revient à la}
\struck{\ill{}} donnée \struck{\ill{}}, pour tt voisinage infinitésimal
$Y'$ d'un $Y \xrightarrow{f} X$, (sur $S$), \add{soumis à la condition
déjà énoncée (\uncertain{automatiquement} si $X/S$ formellement
lisse)}\note{en interligne, en petits caractères} d'un
\emph{prolongement} \struck{\ill{}} de $\underline{F}$ en un
\emph{Module} sur $Y'$, --- avec des conditions de fonctorialité
évidentes. Pour exprimer plus commodément ce point de vue, et ses
variantes, on va introduire un site ad-hoc.

\begin{tikzcd}
X \arrow[dd, no head] & Y \arrow[l, "f"'] \arrow[d, hook] \\
 & Y' \arrow[ul, "f^{\prime}"] \\
S &
\end{tikzcd}

\note{croquis dans la marge gauche, en regard du texte ; les lettres
$Y$ et $Y'$ y sont récrites en encre plus foncée, et l'inclusion
$Y \subset Y'$ est notée par un simple signe $\subset$ vertical. Dans le
texte, « $f : Y \to X$ » est récrit sur une autre lettre}

\page{54}

4.2. Définition du site et du topos \emph{cristallins} de $X$ sur $S$.
Définition du site et du topos stratifiant de $X$ sur $S$.

C'est pareil si $X$ formell\add{ement} lisse sur $S$.

\add{Cas des topos annelés.}\note{en interligne, au-dessus d'une ligne
biffée} \struck{Description d'un}\ Attention, le site n'a \emph{pas de
générateur d'objet final}\,!\note{ce « Attention\ldots{} » est lui-même
écrit en interligne, au-dessus du « Description d'un » biffé} Cas des
\struck{\ill{}} et topos \emph{annelés}.

Description \uncertain{dès} \ill{} faisceaux d'ensembles en termes de
faisceaux habituels. \uncertain{Item} pour faisceaux d'anneaux, faisceaux
de Modules sur faisceaux d'anneaux etc.

Faisceaux spéciaux : les hom.\ de changement de base sont des isom.

Faisceaux de Modules spéciaux sur site stratifiant de $X/S$
$\simeq$\note{le signe est écrit verticalement entre les deux lignes}
Faisceaux de Modules : stratification sur $X/S$.

4.3. Énonçons de suite \uncertain{ces} résultats et conjectures suivants,
qui seront \struck{étudiés} détaillés plus tard :

\emph{Th.} Si $S$ de car.\ $0$ et $X$ lisse sur $S$, alors
\[
  H^{*}_{\substack{\text{cris}/S \\ \text{ou strat}}}(X/S, \mathcal{O}_X)
  \simeq H^{*}_{\mathrm{DR}}(X/S)
\]
\note{sous $H^{*}$, en petits caractères, « cris/$S$ ou strat », en partie
biffé et douteux}
\uncertain{lisse} \uncertain{car} :
\struck{$Rf_{*}(\mathcal{O}_X) \simeq Rf_{*}(\Omega^{\cdot}_{X/S})$ \ldots{}
$(R f_{*}(\mathcal{O}_X))^{\mathrm{cub}} \simeq
Rf_{*}(\Omega^{\cdot}_{X/S}) \ldots$}\note{deux lignes encadrées et
biffées de traits obliques serrés ; lecture très incertaine}

\page{55}

\add{Ceci donnera une explication}\note{en tête de page, au-dessus de la
ligne suivante}

\emph{Conj.} Si $X$ loc.\ de t.f.\ sur $\mathbb{C}$, alors
\[
  H^{*}_{\substack{\text{strat} \\ \text{ou crist.}}}(X/\mathbb{C},
  \mathcal{O}_X) \simeq H^{*}(X^{\mathrm{an}}, \mathbb{C})
\]
\note{l'indice « strat » est récrit, « ou crist.\ » est ajouté au-dessous}

N.B. La conjecture se prouverait, grâce au th., \struck{lorsque}, et au
th.\ de comparaison de l'exposé $1$, dans le cas où $X$ est lisse sur
$\mathbb{C}$. \emph{Donc dans le cas lisse, on trouve d'une façon de
car.\ $0$, une nouvelle description de la coh.\ de D.R.\ sans utiliser
de différentielles.}\note{la phrase est soulignée ligne à ligne ; « une
nouvelle » est récrit, en partie biffé}

\page{56}

\marginal{6.\,12.\,1966}\note{date en haut à droite, dans un cadre ouvert ;
le troisième chiffre est mal formé, lu 6. Elle s'accorde avec le
« 15.11.66 » de la p.\ 35 (batch 2)}

\emph{Laïus}\note{titre souligné, en haut à gauche}

$X \xrightarrow{f} S$\note{écrit verticalement dans la marge gauche :
$X$ au-dessus de $S$, la flèche $f$ descendante}
Soit $\mathcal{F}$ un Module quasi-cohérent stratifié sur $X/S$. On a vu
qu'il définit un module de $X_{\mathrm{strat}}$, et par là un faisceau
cosimplicial \add{$\mathcal{F}^{*} = C^{*}(\mathcal{F})$}\note{entouré, en
interligne} sur $X_{\mathrm{zar}}$, par
\[
  C^{n}(\mathcal{F}) = \mathcal{F}^{(n)} =
  \varprojlim_{i} \mathcal{F}_{\Delta^{n}_{X/S}(i)}
\]
[par n'importe laquelle des $n$ \uncertain{projections}
$\Delta^{n}_{X/S}(i) \to X$ on peut définir l'image inverse, il y a des
isom.\ canoniques entre eux, précisément à cause de la stratification de
$\mathcal{F}$]. \uncertain{On a}
\[
  \boxed{\;H^{*}(X_{\mathrm{str}}, \mathcal{F}) \simeq
  \mathbb{H}^{*}\bigl(X_{\mathrm{zar}}, C^{*}(\mathcal{F})\bigr)\;}
\]
On voit donc deux théories cohomologiques ; \struck{\uncertain{par}}
\add{\uncertain{données}}\note{au-dessus du mot biffé} \struck{\ill{} \ill{}
semble \ill{} \ill{} à première vue \ill{} distinctes \ill{} pas contenir
l'autre}\note{deux lignes biffées}:

a) Pour un $\mathcal{F}$ quasi-cohérent sur $X$, les
$H^{i}(X_{\mathrm{zar}}, \mathcal{F})$.

b) Si $\mathcal{F}$ est \emph{stratifié}, on forme les
$H^{i}(X_{\mathrm{strat}}, \mathcal{F})$, qui s'explicitent par la formule
plus haut.

À première vue, le lien entre ces deux théories cohomologiques semble
donné exclusivement par

\page{57}

la formule ci-dessus ; \struck{Il ne semble pas}
\struck{\uncertain{possible} que} \add{À priori}\note{au-dessus des mots
biffés}, la cohomologie stratifiante d'un Module stratifié s'explicite
comme une hypercohomologie zariskienne (et encore, dans un complexe
dont les \struck{\ill{}} différentielles\note{récrit sur « différentiels »}
n'est pas linéaire\,!), \struck{et} mais l'inverse ne semble pas vrai.
On va voir pourtant que, dans un certain sens, il en est bien ainsi. À
peu de choses près, on associe à tt Module quasi-cohérent $\mathcal{F}$
(\emph{sans stratification}) un Module \struck{quasi-cohérent}
\add{(plutôt, un pro-Module pro-quasi-cohérent)}\note{en interligne,
relié par un trait} $\mathcal{Q}^{\cdot}(\mathcal{F})$ stratifié, dont la
cohomologie stratifiante est la coh.\ zariskienne \emph{ordinaire} de
$X$. De plus, $\mathcal{Q}^{\cdot}(\mathcal{F})$ est un foncteur en
$\mathcal{F}$, non seulement par des homomorphismes linéaires
$\mathcal{F} \to \mathcal{G}$, mais $=$ \struck{\ill{}} \add{(\uncertain{conoyaux}}
\struck{additifs, \ill{} \ill{}} de $\mathcal{O}_X$-modules) par des
\emph{D} : $\mathcal{F} \to \mathcal{G}$ qui sont seulement des op.\
diff.\ rel/$S$.\note{le passage « On va voir\ldots{} rel/$S$ » est marqué
d'un trait vertical dans la marge gauche ; « conoyaux » est écrit au-dessus
de la rature}
\struck{\ill{}}\ \add{[Dans le cas très}\note{au-dessus d'une ligne
biffée} particulier où $\mathrm{diag}_{X/S} : X \to X \times_S X$ est une
immersion nilpotente, donc \struck{\ill{} $\mathcal{Q}^{\cdot}(\ill{})$}
stratification $\equiv$ donnée de descente : $S$,
$\mathcal{Q}^{\cdot}(\mathcal{F}) = f^{*}f_{*}(\mathcal{F})$].
L'intérêt \add{de ces\ldots{} des \uncertain{uns voisins}}\note{ajout
entouré, en interligne, en partie illisible} \struck{\ill{}} est que
l'\emph{hypercoh.}\ \add{zariskienne ordinaire}\note{au-dessus de la
ligne} \struck{d'un} \emph{complexe d'op.\ diff.}\ \emph{peut
s'interpréter comme la coh.\ stratifiante des complexes commutant de
faisceaux stratifiés\ldots\,!}\note{les deux dernières lignes sont
serrées en bas de page ; la dernière est d'une autre encre, plus pâle}

\page{58}

\ill{} les \emph{op.\ diff.}\ de $M$ dans $N$ (rel.\ à $S$), alors on
trouve encore un foncteur \add{\uncertain{exact}}\note{au-dessus de la
ligne} \struck{\ill{}} $M$, mais non à valeurs dans les
$\mathcal{P}^{i}$-Modules, \add{(si $i \geq 1$)}\note{encadré, en
interligne} mais seulement dans les $\mathcal{P}^{i-1}$-Modules, si on
considère $\mathcal{P}^{i}[M]$ comme $\mathcal{P}^{i-1}$-Module par
restriction des scalaires \add{(On les note alors
$\mathcal{Q}^{i-1}[M]$)}\note{en interligne, avec un renvoi marqué d'un
astérisque cerclé vers la formule encadrée du bas de page} à l'aide de
$\alpha^{i-1} : \mathcal{P}^{i-1} \to \mathcal{P}^{i}$. En particulier,
$\mathcal{Q}^{0}[M]$ est \struck{\ill{}} un foncteur, en
$M \in \mathrm{Ob}\,\mathrm{Mod}(X/S)$ \add{$\to$ [Modules sur
$\mathcal{O}_X$, avec \uncertain{les} op.\ diff.\ \struck{d'ordre}
rel.\ à $S$],}\note{ajout relié par une flèche à « $\mathrm{Mod}(X/S)$ »}
\ill{} \ill{} à valeurs dans les pro-$\mathcal{O}_X$-Modules.
\marginal{\ill{} \ill{} \ill{} \ill{} \struck{\ill{}} \ill{} \ill{}
\emph{Anti-Riesz} \ill{} \ill{} \ill{} \ill{} $X$ \ill{}
\uncertain{pro-objet} \ill{} \ill{}]}\note{note écrite verticalement dans
la marge gauche, en quatre ou cinq lignes, reliée au texte par un trait
terminé par un petit cercle ; seul « Anti-Riesz » se lit avec quelque
sûreté}

\struck{[$\mathcal{P}^{0}[M] = M$ \ill{} \ill{} \ill{} \ill{} \ill{}
$\mathrm{Ab}(X)$]}\note{deux lignes entre crochets, encadrées et barrées
de traits obliques}

Notons que les \ill{} $\alpha^{\cdot} : \mathcal{P}^{*} \to
(S\,\mathcal{P})^{*} = \mathcal{Q}^{*}$ permettent de considérer
$\mathcal{Q}^{*}$ comme \add{algèbre}\note{au-dessus d'un mot biffé}
\struck{\ill{}} cosimplicial\add{e}\note{« cosimplicial » est récrit sur
un autre mot} sur l'\struck{\ill{}} \uncertain{anneau} cosimplicial
$\mathcal{P}^{*}$, et que les opérations extérieures \uncertain{ducales}
\[
  \boxed{\;(*)\quad \mathcal{Q}^{i}[M] = \mathcal{P}^{i}
  \otimes_{\mathcal{O}_X} M\;}
\]
\note{la formule, en bas de page, est encadrée et marquée d'un astérisque
cerclé ; sur la page l'indice de $\mathcal{Q}$ se lit « $i{-}1$ » ou
« $i$ » ; le mot avant la formule, « ducales », est très douteux}

\page{59}\note{plan de l'exposé, en tête de page, séparé de la suite par
un trait horizontal}

\begin{itemize}
  \item[§\,1] Coh.\ de De Rham ordinaire
  \item[§\,2] Tour d'horizon sur coh.\ $\ell$-adiques. Motivations pour
  le désir d'une coh.\ $p$-adique en car.\ $p$.
  \item[§\,3] \uncertain{Topo} de Washnitzer-Monsky.
  \item[§\,4] Stratifications. Relation de W-M.\ avec connexions de
  Gauss-Manin sur coh.\ de De Rham
  \item[§\,5] \struck{Le} Site cristallin et site stratifiant. Énoncé
  \struck{\ill{}} d'un théorème et d'une conjecture en car.\ nulle.
  \item[§\,6] Calculs Čechistes (car.\ quelconque).
  \item[§\,7] Comparaison de la Coh.\ de De Rham et cristalline (car.\
  \struck{nulle}).\note{« nulle » est barré d'un trait}
  \item[§\,8] \struck{\uncertain{Comparaison}} \add{Retour}\note{au-dessus
  du mot biffé} sur le cas de la car.\ $p$\ldots
\end{itemize}

\marginal{\struck{\ill{}} Coh.\ stratifiante}\note{dans la marge gauche,
en face de 6.1, soulignée ; un premier titre y est noirci}
6.1. Commençons à travailler sur le \struck{\ill{}} \emph{topos}
\emph{stratifiant} de $X/S$. Il y a l'objet $\widetilde{X}$ de ce topos :
il \emph{couvre l'objet final} (c'est ici qu'on utilise la restriction sur
ces objets du site stratifiant parmi ceux du site cristallin). Donc suite
spectrale de Leray
\[
  H^{*}(X_{\mathrm{strat}}, F) \Longleftarrow
  E_2^{pq} = H^{p}\bigl(\nu \mapsto H^{q}(\widetilde{X}^{\nu+1}, F)\bigr)
\]
On a aussi que
\[
  \widetilde{X}^{\nu+1} = \varinjlim_{i} \Delta^{(\nu)}_{X}(i)
\]
\note{une rature noircit l'exposant de $\widetilde{X}$ (récrit
« $\nu+1$ ») et un symbole à droite de $\Delta^{(\nu)}_X(i)$ ; une
accolade relie ce dernier à une glose}
[$\Delta^{(\nu)}_X(i)$ : voisinage inf.\ $i$-ième de la diagonale dans
$X^{\nu+1}$, considéré comme \uncertain{prosimplicial} \uncertain{au-dessus}
de $X$ ; $X \hookleftarrow U \to X'$, $(U, X')$
\struck{\ill{}}]\note{la glose est écrite sous la formule, en petits
caractères, avec un croquis $X \leftarrow U \to X'$}

D'autre part, pour tt objet du site stratifiant, on a

\page{60}

\[
  H^{*}\bigl((U, X'), F\bigr) \simeq H^{*}(X', F_{(U, X')})
\]
Supposons pour simplifier $X$ affine, \struck{et affine sur $S$,} donc
les $\Delta^{(\nu)}(i)$ affines. Supposons que les $F_{\Delta^{(\nu)}(i)}$
soient \emph{quasi-cohérents} \add{et \uncertain{indépendants}
\uncertain{mutuellement}}\note{en interligne, souligné, relié par un
trait}. Alors on trouve donc
\[
  H^{q}\bigl(\widetilde{\Delta^{(\nu)}_X(i)}, F\bigr) = 0
  \quad \text{si } q > 0
\]
$H^{0}\bigl(\widetilde{\Delta^{(\nu)}_X(i)}, F\bigr)$ pour $i$ variable
\uncertain{est un syst.\ proj.\ strict}, d'où par les arguments : la
M.L.\note{sans doute la condition de Mittag-Leffler}
\[
  H^{q}(\widetilde{X}^{(\nu+1)}, F) =
  H^{q}\Bigl(\varinjlim_{i} \widetilde{\Delta^{(\nu)}_X(i)}, F\Bigr) =
  \varprojlim H^{q}\bigl(\widetilde{\Delta^{(\nu)}_X(i)}, F\bigr)
\]
donc
\[
  \left\lbrace
  \begin{array}{l}
    H^{q}(\widetilde{X}^{(\nu+1)}, F) = 0 \quad \text{si } q > 0 \\
    H^{0}(\widetilde{X}^{(\nu+1)}, F) = \varprojlim_{i}
    F\bigl(\Delta^{(\nu)}_X(i)\bigr)
  \end{array}
  \right.
\]
et par suite des isom.\ can.
\[
  \boxed{\;H^{*}(X_{\mathrm{strat}}, F) \simeq
  H^{*}\bigl(\nu \mapsto F(X^{(\nu+1)}/X)\bigr)\;}
\]
\note{sous $X^{(\nu+1)}/X$, une accolade : « complété formel » ; le
dénominateur est récrit sur une rature}
si $X$ affine, les $F_{\Delta^{(\nu)}_X(i)}$ quasi-cohérents
(p.\,ex.\ $F = \mathcal{O}_{X_{\mathrm{strat}}}$).

Dans le cas où on \struck{\ill{}} \add{\uncertain{suppose}}
\uncertain{seulement} que les $F_{\Delta^{(\nu)}_X(i)}$ sont
qu.-coh.\ \add{et \uncertain{indépendants} mutuellement}\note{ajout en
interligne, au-dessus d'un passage biffé et en partie illisible},
mais $X$ pas nécessairement affine, on \uncertain{montre} que pour
tout entier $\nu$, \ill{} a une

\end{document}
